% Part: computability
% Chapter: recursive-functions
% Section: trees

\documentclass[../../../include/open-logic-section]{subfiles}

\begin{document}

\olfileid{cmp}{rec}{tre}
\olsection{树}

有时将树表示为自然数会很有用，
就像我们可以用数来表示序列，并通过针对其编码的原始递归关系与函数来处理序列的性质与运算一样。
我们将利用序列及其编码来实现。
一棵树可以是一个单独的节点（可能带有标签），
也可以是一个连接着若干子树的节点（可能带有标签）。
这个节点被称为该树的\emph{根}，它所连接的子树称为其\emph{直接子树}。

我们递归地将树编码为一个序列 $\tuple{k, d_1, \dots, d_k}$，
其中 $k$ 是直接子树的个数，而 $d_1$, \dots,~$d_k$
是这些直接子树的编码。如果节点带有标签，标签可以放在直接子树编码之后。
因此，仅由带标签~$l$ 的单个节点构成的树编码为 $\tuple{0,l}$，
而由根（带标签~$l_1$）连接两个单节点（分别带标签 $l_2$ 和 $l_3$）构成的树编码为 $\tuple{2,
  \tuple{0, l_2}, \tuple{0, l_3}, l_1}$。

\begin{prop}
  \ollabel{prop:subtreeseq}
  函数 $\fn{SubtreeSeq}(t)$ 返回一个序列的编码，该序列的元素为编码为~$t$ 的树的所有子树的编码。此函数是原始递归的。
\end{prop}

\begin{proof}
  首先注意 $\fn{ISubtrees}(t) = \fn{subseq}(t, 1, (t)_0)$ 是原始递归的，它返回树~$t$ 的直接子树的编码。
  现在我们可以定义一个辅助函数 $\fn{hSubtreeSeq}(t,n)$，
  用来计算所有与根节点相距 $n$ 个节点的子树组成的序列。
  与~$t$ 的根节点距离为 $0$ 的子树序列——换言之，以~$t$ 的根节点为起点的子树——就是仅由~$t$ 构成的序列。
  要获得~$t$ 的所有第 $n+1$ 层子树的序列，我们将第 $n$ 层子树的序列与一个由第 $n$ 层子树的直接子树组成的序列拼接起来。
  要得到所有这些子树的列表，注意：如果 $f(x)$ 是返回序列编码的原始递归函数，那么 $g_f(s, k) = f((s)_0) \concat
  \dots \concat f((s)_k)$ 也是原始递归的：
    \begin{align*}
      g(s, 0) & = f((s)_0)\\
      g(s, k+1) & = g(s, k) \concat f((s)_{k+1})
    \end{align*}
    例如，如果 $s$ 是一个树的序列，那么
    $h(s) = g_{\fn{ISubtrees}}(s, \len{s})$ 给出~$s$ 中元素的直接子树的序列。
    我们可以借此定义 $\fn{hSubtreeSeq}$ 如下：
    \begin{align*}
      \fn{hSubtreeSeq}(t, 0) & = \tuple{t} \\
      \fn{hSubtreeSeq}(t, n+1) & = \fn{hSubtreeSeq}(t, n) \concat
      h(\fn{hSubtreeSeq}(t, n)).
    \end{align*}
    由~$t$ 编码的树中子树的最大层数，即根节点与叶节点之间的最大距离，以编码~$t$ 为上界。
    因此，由编码~$t$ 表示的树的所有子树的编码序列由 $\fn{hSubtreeSeq}(t, t)$ 给出。
\end{proof}

\begin{prob}
  \olref[cmp][rec][tre]{prop:subtreeseq} 的证明中对 $\fn{hSubtreeSeq}$ 的定义一般会包含重复项。
  请给出另一种定义，保证每个子树的编码在结果列表中只出现一次。
\end{prob}

\end{document}